\clearpage
% ====================================================================
% Section: Conclusions and Future Work
% ====================================================================
\chapter{Conclusions and Future Work}\label{cap:conclusion}

This dissertation has presented a kill-chain-aware adaptive defense framework for IoT networks that couples a finite-budget, reactive tug-of-war attacker, supervised stage detection, and Deep Reinforcement Learning under a unified evaluation contract in which the defender acts on raw feature rows without ever observing the true kill-chain stage. The framework was developed and rigorously evaluated on the contemporary CICIoT2023 dataset across an audit-first, eight-stage development protocol. This concluding chapter summarises the contributions (Section~\ref{sec:summary}), discusses the empirical findings worth defending (Section~\ref{sec:findings}), surfaces the limitations of the approach (Section~\ref{sec:limitations}), and outlines a roadmap of future research directions (Section~\ref{sec:future_work}).

% --- Subsection 5.1 ---
\section{Summary of Contributions}\label{sec:summary}
The proliferation of IoT devices has created a vast and vulnerable attack surface, overwhelming conventional security paradigms that rely on static signatures and manual intervention. This dissertation has addressed this critical gap with a closed-loop security framework that synergises predictive analytics with autonomous, adaptive defense. The primary contributions are:
\begin{enumerate}
    \item \textbf{A kill-chain-aware adaptive defense architecture under partial observability.} A two-stage framework integrating a \textbf{tug-of-war attacker}, operating under a finite intrusion budget, and a supervised stage detector with a five-action DRL defender (\textsc{observe}, \textsc{log}, \textsc{restrict}, \textsc{block}, \textsc{isolate}) trained against a stage-action proportionality reward. At each active stage the attacker reacts to the defender's force: a proportionate action pushes it one stage back the kill chain, an under-forced action lets it advance, and an over-forced action merely holds it in place. Critically, the defender observes only the 29-dimensional per-flow feature vector and \textbf{never the true attack stage}, so the control task is fundamentally one of stage inference under partial observability (a POMDP); this perception gap, rather than control complexity, is the central difficulty the learned policy must overcome.

    \item \textbf{Empirical grounding on a contemporary IoT dataset, with an honest leakage-bug fix.} The framework is developed and validated on the \textbf{CICIoT2023 dataset}, projected onto a five-stage Cyber Kill Chain. A subtle splits-construction bug --- which had inadvertently leaked 70\,\% of every held-out out-of-distribution attack class back into the training pool --- was discovered during stage-detector training, fixed at commit \texttt{3cd2fb9}, and locked behind disjointness assertions before any benchmark numbers were collected. The fix and its forensics are presented honestly in Section~\ref{sec:dataset_results}.

    \item \textbf{A comparative DRL benchmark framed as an honest deployment trade-off.} DQN, PPO, and A2C were trained over \NumSeeds{} seeds and evaluated on a held-out balanced test split ($n=300$ episodes per policy) against four reference policies (random, always-\textsc{observe}, always-\textsc{block}, RF-Acting) and an oracle Recommended-Action rule that observes the true attack stage. The three trained agents are statistically indistinguishable on mean reward (overlapping bootstrap confidence intervals), and dominate every non-RL \emph{trivial} baseline; the best deployable agent (\BestAgentName, $+\BestAgentReward$) captures \textbf{\OracleCapturePct\,\%} of the oracle ceiling ($+\OracleCeiling$) without ever observing the true attack stage, while keeping its benign false-positive rate below 1\,\% (all three agents fall in the \FPRdqnPct--\FPRppoPct\,\% range). We do not claim DRL dominates the strongest deployable baseline: a detector-coupled Random-Forest policy (RF-Acting) attains a higher reward ($+\RFReward$) but depends on an upstream stage classifier and incurs a \RFLatency\,ms per-decision latency, whereas the trained RL policy is detector-free at $\sim$\BestAgentLatency\,ms. A measured latency study (F7b) shows this gap is an artefact of single-sample classifier dispatch rather than a fundamental cost; the genuine RL advantage is detector-independence and a fixed sub-0.1\,ms decision with no upstream model in the loop.


    \item \textbf{Reward mis-specification: discovered, diagnosed, and structurally fixed.} A 12-cell sparse one-at-a-time reward-component sweep identified that the reward-maximising policy under the naive reward contract ($\texttt{impact\_is\_terminal}=\mathit{True}$) diverges from the human-intended objective (``defend the \textsc{impact} step'') because de-escalation bonuses dominate the \textsc{impact}-step penalty. This is reward mis-specification in the formal sense~\cite{Amodei2016,Krakovna2020}. A \emph{structural} environment change --- setting $\texttt{impact\_is\_terminal}=\mathit{False}$ (the corrected primary contract) --- resolves the mis-specification at its source. The structural-fix ablation strand reports a security-KPI mitigated-impact rate of $\FnineStructuralMitRate$ for the structural fix, versus $0.0$ for the mis-specified baseline contract under the same strand --- a decisive improvement on the meaningful security objective. Complementing the structural fix, a \emph{finite intrusion budget} ($B=40$, locked) makes prevention --- exhausting the attacker before \textsc{impact} --- a first-class, measurable outcome: under the primary contract \emph{prevention} becomes the defender-attributable security KPI, with the stage-aware oracle preventing every episode (prevention rate $1.00$) at near-zero availability cost while the trained agents prevent between $0.33$ and $0.60$ of episodes. A six-point sweep over the environment-difficulty parameter $p_{\mathrm{down}}$ (the tug-of-war de-escalation success probability, headline $0.90$) characterises the trained policy's reward as monotone non-decreasing in environment leniency. Evaluation against four held-out OOD attack classes yields the sharpest single demonstration of detector independence: on \texttt{VulnerabilityScan}, the class on which the supervised stage detector is structurally blind (recall $0.001$), the detector-free trained policy ($+298.3$, within seed-noise of its in-distribution mean) \emph{decisively outperforms} the detector-coupled RF-Acting baseline ($-4430.6$), which collapses because its blind classifier systematically under-forces against an advancing attacker; the gap is $\Delta=+4728.9$ and gate~G7.9 passes.

    \item \textbf{An audit-first reproducibility protocol.} Every figure ships a \texttt{manifest.json} that pins the SHA-256 hashes of every input artefact and the producing Git commit. A smoke-reproducibility harness (\texttt{python -m scripts.reproducibility\_smoke}) verifies the full hash chain end-to-end on a fresh checkout in approximately five seconds, returning a verdict bucketed into \textsc{ok} / \textsc{fail} / \textsc{known-divergence} / \textsc{skip}. The complete protocol is documented in Appendix~A.
\end{enumerate}

% --- Subsection 5.2 ---
\section{Principal Empirical Findings}\label{sec:findings}
Three findings are pre-registered, principled results of the approach rather than ad-hoc failures, and are arguably the strongest contributions of this dissertation:

\begin{itemize}
    \item \textbf{Finding 1 (Blue-Team Training stage, gate G5.4 \textsc{pass-with-finding}): Reward Mis-specification, Discovered and Structurally Fixed.} A reward-decomposition trace explains the gap between high episodic reward and low mitigated-impact rate as principled reward mis-specification~\cite{Amodei2016,Krakovna2020} under the naive contract ($\texttt{impact\_is\_terminal}=\mathit{True}$): de-escalation bonuses earned during the kill chain dominate the expected \textsc{impact}-step loss, so the reward-maximising policy --- ``rack up de-escalations and accept the \textsc{impact} loss'' --- is \emph{correct} under the proxy objective but \emph{wrong} under the intended objective. The finding is not a regression: it is the input that motivates the reward-component ablation, which identifies and applies the structural fix ($\texttt{impact\_is\_terminal}=\mathit{False}$, the corrected primary contract). The structural-fix ablation strand reports a security-KPI mitigated-impact rate of \FnineStructuralMitRate, versus $0.0$ for the mis-specified baseline contract under the same strand. Notably, the redesigned reward (capped proportionality and de-escalation bonuses) also removes the over-blocking incentive that previously inflated benign false positives: all three trained agents now keep their benign FPR below 1\,\% (\FPRdqnPct--\FPRppoPct\,\%), so the structural fix improves the security KPI without the operational FPR cost that an earlier reward design had incurred.

    \item \textbf{Finding 2 (Held-Out Benchmark stage, audit AF2 reframe of gate G6.2): \OracleCapturePct\,\% of the Oracle Ceiling.} The best trained DRL agent (\BestAgentName, $+\BestAgentReward$) captures \OracleCapturePct\,\% of the oracle ceiling ($+\OracleCeiling$ on \texttt{test\_balanced}) without observing the true attack stage; the oracle is reported as a measurement instrument rather than a competing baseline because it reads the true attack stage directly from the simulator and cannot be deployed. The AF2 reframe is preserved in the JSON scoreboard record (the original ``RL $>$ recommended-action rule'' threshold reads \texttt{passes: false} permanently), so it is not goalpost-moving --- it is an interpretation correction documented in the results prose. All three RL agents are statistically indistinguishable from each other on mean reward (\NumSeeds{} seeds, overlapping CIs); the main deployable trade-off is between RF-Acting (higher reward, but detector-dependent and \RFLatency\,ms p50 latency) and trained RL (lower reward, detector-free, \BestAgentLatency\,ms latency). A measured latency study (F7b) shows the \LatencyRatio$\times$ gap is an artefact of single-sample classifier dispatch rather than a fundamental cost of supervised models; the genuine, defensible RL advantage is detector-independence --- a fixed sub-0.1\,ms decision with no upstream classifier in the loop.

    \item \textbf{Finding 3 (Ablation \& Robustness stage, gates G7.2 \textsc{pass-with-finding} and G7.9 \textsc{pass}): Structural Fix and OOD Detector-Independence Dividend.} A 12-cell reward-coefficient sweep characterises the limit of one-at-a-time coefficient shaping: 11 of 12 cells stay within $\pm 150$ reward of the centre baseline. The structural fix ($\texttt{impact\_is\_terminal}=\mathit{False}$) achieves an ablation-strand mean reward of $+\FnineStructuralReward$ and mitigated-impact rate \FnineStructuralMitRate. Out-of-distribution evaluation against the hardest held-out class \texttt{VulnerabilityScan} (supervised-detector recall $0.001$) makes the detector-independence argument concrete: the trained DRL policy (\BestAgentName{} $+298.3$) holds near its in-distribution mean, while the detector-coupled RF-Acting baseline collapses to $-4430.6$ because its blind classifier systematically under-forces against an advancing attacker. Gate~G7.9 passes ($\Delta=+4728.9$ in the learned policy's favour), and the result is promoted to a Tier-1 deliverable (audit finding AF1): the supervised baseline's higher in-distribution reward comes bundled with a catastrophic out-of-distribution failure mode that the detector-free policy avoids.
\end{itemize}

These three findings, together with the audit-first reproducibility protocol that makes them verifiable, constitute the empirical backbone of the dissertation.

% --- Subsection 5.3 ---
\section{Limitations}\label{sec:limitations}
The work has four principal limitations, surfaced honestly in Section~\ref{sec:discussion} and recapitulated here.

\begin{enumerate}
    \item \textbf{Prevention depends on the finite intrusion budget.} Without a budget, an attacker that the defender fails to suppress eventually advances every episode to \textsc{impact}. The finite intrusion budget ($B=40$, locked) changes this: by draining the attacker on every active progression step and on every defender-forced de-escalation, it makes \emph{prevention} --- exhausting the attacker before \textsc{impact} --- the primary, defender-attributable security KPI (replacing the degenerate \texttt{mitigated\_impact\_rate}, which collapses onto \texttt{compromise\_rate} for fixed always-\textsc{block} policies). Prevention is then a policy-dependent outcome: the stage-aware oracle prevents every episode (prevention rate $1.00$) at near-zero availability cost, the trained agents prevent between $0.33$ and $0.60$, and always-\textsc{block} prevents every episode only by indiscriminate force (full availability cost), making it the worst-scoring policy on reward. The budget-sweep ablation (F16) confirms this prevention peaks at an interior budget and vanishes at both extremes, so the headline prevention result is contingent on this designed budget mechanism; under an unbounded attacker the trained policies remain \emph{reactive mitigation} agents.

    \item \textbf{MTTC values bounded by the lifecycle-floor \textsc{impact} clamp.} The clamp downgrades pre-step-20 \textsc{impact} transitions to \textsc{maneuver}, so mean MTTC (observed in the $\sim$24--25-step range for trained agents) is biased upward by the $\texttt{min\_episode\_length} = 20$ floor. Mean MTTC numbers should not be interpreted as ``unconstrained natural'' MTTC.

    \item \textbf{Trained RL does not beat RF-Acting on the in-distribution \texttt{test\_balanced} split.} The deployable comparison is a deployment trade-off, not a reward-domination claim: on the in-distribution test split RF-Acting achieves higher reward ($+\RFReward$) but couples to an upstream stage classifier and incurs \RFLatency\,ms p50 inference latency, whereas the trained policy is detector-free at \BestAgentLatency\,ms. The measured latency study (F7b) shows the \LatencyRatio$\times$ gap is an artefact of single-sample classifier dispatch rather than a fundamental cost; the dissertation frames the trade-off honestly and does not claim RL dominates RF-Acting in absolute in-distribution reward. This in-distribution edge does not extend to distribution shift: on the hardest OOD class (\texttt{VulnerabilityScan}) the ordering reverses entirely, with the detector-free policy outperforming RF-Acting by $\Delta=+4728.9$ (Finding~3).

    \item \textbf{Designed (not learned) attacker dynamics.} The tug-of-war attacker reacts to the defender's force --- a proportionate response pushes it one stage back, an under-forced response lets it advance, and an over-forced response merely holds it --- but it is not jointly or adversarially co-trained with the defender, and its budget mechanics are designed rather than learned, limiting realism relative to a real adversary that may pivot, retreat, interleave kill-chain stages, or adapt to the defender over time. This is a deliberate, explicitly-bounded simplification; the F17 evasion-reactive sweep (gate G7.10) shows the learned defender degrades gracefully when the attacker stalls in response to defensive force, but a fully adaptive, co-adapting, or non-monotonic attacker is deferred to future work.
\end{enumerate}

% --- Subsection 5.4 ---
\section{Future Work and Research Directions}\label{sec:future_work}
The audit-first protocol that produced this dissertation explicitly delineated work scoped \emph{into} this dissertation from work scoped \emph{out of} it. Six follow-up directions are surfaced here in order of mentor-recommended priority.

\subsection{Tightly Scoped Extensions of Empirical Findings}

\paragraph{Direction 1 --- Non-monotonic and adaptive attacker model.} The current $5\times5$ tug-of-war attacker is upper-triangular in its autonomous dynamics: absent defensive force the kill chain advances but never retreats on its own (defender-driven de-escalation already pushes it back down one stage under a proportionate action). Relaxing this to a non-monotonic attacker (a fuller transition matrix, or a learned sequence model) that allows stage pivots and retreats would produce more realistic episode trajectories and test the defender's robustness to adversarial pivoting. This is the most direct extension of the monotonic-attacker limitation, as the training environment's reward function already supports de-escalation; only the attacker's transition dynamics need to change. The F17 evasion-reactive coupling is a first step in this direction --- the attacker already stalls in response to defensive force --- and a fully adaptive, self-play attacker is the natural continuation (Direction 4).

\paragraph{Direction 2 --- Edge-hardware deployment and latency budget.} The \LatencyRatio$\times$ latency advantage of trained \BestAgentName{} over RF-Acting (\BestAgentLatency\,ms vs.\ \RFLatency\,ms p50) establishes trained RL as the candidate policy for resource-constrained edge-hardware deployments. A natural follow-up is to quantify the latency and accuracy trade-off on actual IoT gateway hardware (e.g.\ ARM Cortex-A-class processors) under real packet rates, and to apply model-compression techniques (quantisation, pruning) to the trained policy network to fit within the tighter computational budgets of embedded defenders. The policy's $\sim$4K-parameter footprint ($\sim$20\,KB) already makes it competitive with embedded ML models.

\paragraph{Direction 3 --- Train-time OOD-class augmentation.} The F15 evaluation establishes detector-independence on a single held-out class; the natural next step is to test how far that robustness generalises. Instead of evaluating against held-out classes one at a time, \emph{include} simulacra of a broader OOD family in training (synthetic feature blending, domain randomisation, or a small adversarial-augmentation generator) and measure whether the detector-free advantage observed on \texttt{VulnerabilityScan} extends to harder and more diverse distribution shifts. This probes whether the catastrophic-failure-avoidance demonstrated on one class is a general property of detector-free control or a class-specific artefact.

\subsection{Broader Research Directions}

\paragraph{Direction 4 --- Multi-agent reinforcement learning and self-play adversary.} A real IoT deployment includes many overlapping defenders (per-segment defenders coordinating with a central Security Operations Centre). This can be modelled as a cooperative multi-agent RL (MARL) system where specialised agents defend different network segments and learn to coordinate. Extending to a self-play adversary --- where a learned attacker is jointly trained with the defender to produce adversarial stage trajectories, replacing the fixed Markov transition matrix --- would produce a harder, more dynamic training distribution that could close the remaining gap to the oracle ceiling without a structural reward-function change. Both extensions inherit the kill-chain decision frame and the audit-first reproducibility protocol from this dissertation.

\paragraph{Direction 5 --- Federated learning for multi-site deployment.} A single-site RL policy trained on one IoT deployment may not generalise to a different network topology or device mix. Federated extensions would allow defender policies to be trained across diverse IoT deployments without centralising sensitive traffic data, producing a population-level policy that is more robust to site-specific distribution shift. This direction is complementary to Direction 4 and is directly motivated by the OOD evaluation's finding that per-site feature novelty limits generalisation.

\paragraph{Direction 6 --- Constrained-MDP formulation for harder operating regimes.} Although the redesigned reward keeps benign-FPR below 1\,\% for all three agents under the headline contract --- the per-episode caps on the proportionality and de-escalation bonuses removed the over-blocking incentive that previously inflated it --- a real deployment may demand a \emph{hard}, certified ceiling on false positives across all regimes (harsher attacker budgets, adaptive adversaries, or distribution shift). A principled extension would express this as a constrained MDP, with the benign-FPR as an explicit constraint enforced via a Lagrange multiplier, or a hierarchical policy that separates the stage-inference decision from the action-selection decision. This would carry the empirical FPR guarantee from a single calibrated regime to a worst-case bound, which is the form an operational safety case ultimately requires.

\subsection{Note on the Environment-Perturbation Robustness Stage}
The audit-first protocol that structured this dissertation defined eight development stages (stages 1--7 delivered; stage 8 reserved for environment-perturbation robustness work including noise-robustness (F13) and OOD-augmentation (F14)). The environment-perturbation robustness stage was \emph{skipped} as a dissertation deliverable: the F9 structural change resolved the reward mis-specification using only the $\texttt{impact\_is\_terminal}$ flag, making a full perturbation sweep less informative for the dissertation's headline claim than a precise articulation of the existing ablation-stage findings. The F13 and F14 directions are reframed above as Directions 3 and 2 (post-thesis future work). This decision is part of the audit-trail anchored at the four committed \texttt{G[N]\_scoreboard.json} files.
